\documentclass[11pt]{article}
\usepackage{lmodern}
\usepackage[T1]{fontenc}
\usepackage[margin=1.05in]{geometry}
\usepackage{amsmath,amssymb,amsthm}
\usepackage[hidelinks]{hyperref}
\newtheorem{theorem}{Theorem}[section]
\newtheorem{corollary}[theorem]{Corollary}
\newtheorem{proposition}[theorem]{Proposition}
\DeclareMathOperator{\Pic}{Pic}
\newcommand{\Z}{\mathbf Z}
\newcommand{\Q}{\mathbf Q}
\newcommand{\C}{\mathbf C}
\newcommand{\R}{\mathbf R}
\newcommand{\F}{\mathbf F}
\title{Congruence kernels of mod-two elliptic period characters}
\author{\small Contact: \href{mailto:admin@cosmosgeek.org}{admin@cosmosgeek.org}}
\date{Working draft}
\begin{document}
\maketitle

\begin{abstract}
Let \(E\) be an optimal elliptic quotient of \(J_0(N)\).
We identify the mod-two elliptic period characters with congruence kernel
with the intersection of \(E^\vee[2]\) and the Shimura subgroup in \(J_0(N)\).
In particular, if \(E(\Q)[2]=0\), all three nonzero period characters have
noncongruence kernel. The argument applies Wohlfahrt's cusp-width theorem and
the description of the Shimura subgroup by Ling and Oesterl\'e. An explicit
formula identifies the primitive imaginary period character in either
real-component case.
\end{abstract}

\section{Period characters}

Let \(q:J_0(N)\twoheadrightarrow E\) be a quotient over \(\Q\) with
connected kernel, and write \(\varphi:X_0(N)\to E\) for the composite with
the Abel--Jacobi map based at \(\infty\). Let \(\pi:X_1(N)\to X_0(N)\)
be the natural map. Using the canonical principal polarization on \(J_0(N)\),
write
\[
 \iota=q^\vee:E^\vee\hookrightarrow J_0(N),\qquad
 \Sigma(N)=\ker\big(\pi^*:J_0(N)\to J_1(N)\big).
\]
Here \(E^\vee=\Pic^0(E)\), and \(\iota\) on line bundles is the map induced
by \(\varphi^*\). We work with complex points unless a ground field is specified.

The lattice map \(\varphi_*:H_1(X_0(N),\Z)\to H_1(E,\Z)\) is onto.
Indeed the Abel--Jacobi map identifies the homology of the curve with that
of its Jacobian. The exact sequence of complex tori for \(q\) then gives
surjectivity on fundamental groups, since the kernel is connected.
For \(\lambda\in\operatorname{Hom}(H_1(E,\Z),\F_2)\), pull back by
\(\varphi_*\) and then evaluate on loops from \(\Gamma_0(N)\). This gives
a character \(\chi_\lambda:\Gamma_0(N)\to\F_2\). One may take the loop
to be the image of a path from \(z\) to \(\gamma z\) in the upper half-plane
for any ordinary base point \(z\). These loops generate the fundamental
group of the compact curve: remove cusps and elliptic points, lift loops
on the complement, and then fill the punctures. Thus \(\chi_\lambda\)
is nonzero whenever \(\lambda\) is nonzero.

Identify \(\F_2\) with \(\mu_2=\{1,-1\}\) by \(a\mapsto(-1)^a\).
The analytic Kummer sequence identifies \(H^1(E,\mu_2)\) with
\(\Pic(E)[2]=E^\vee[2]\), since squaring is surjective on global nonzero
constant functions. Denote the point corresponding to \(\lambda\) by
\(P_\lambda\). Functoriality gives \(\iota(P_\lambda)\) as the two-torsion
line bundle on \(X_0(N)\) for the pulled-back character.

\begin{theorem}\label{th:criterion}
For a nonzero \(\lambda\), put \(H=\ker\chi_\lambda\). The following
conditions are equivalent:
\begin{enumerate}
\item \(H\) is a congruence subgroup of \(\mathrm{SL}_2(\Z)\).
\item \(\Gamma_1(N)\subseteq H\).
\item \(\iota(P_\lambda)\in\Sigma(N)[2]\).
\end{enumerate}
Consequently the group of elliptic period characters with congruence kernel,
including the zero character, identifies with
\(\iota(E^\vee[2])\cap\Sigma(N)[2]\).
\end{theorem}

\begin{proof}
Every cusp parabolic maps to a loop which contracts after the cusp is filled.
The character also kills elliptic stabilizers and \(-I\), since it arises
from the compact curve. Each stabilizing translation at a cusp of
\(\Gamma_0(N)\) therefore has exactly the same minimal width in \(H\).
The least common multiple of these widths is \(N\): the width at the
cusp \(0\) is \(N\), and all widths divide \(N\) since
\(\Gamma(N)\subseteq\Gamma_0(N)\). In particular this is also the least
common multiple of the widths of \(H\). If \(H\) is congruence,
Wohlfahrt's theorem for subgroups containing \(-I\)
\cite[Theorem~2]{Woh} gives \(\Gamma(N)\subseteq H\). Modulo \(\Gamma(N)\)
the group \(\Gamma_1(N)\) is generated by the image of
\(\left(\begin{smallmatrix}1&1\\0&1\end{smallmatrix}\right)\).
This translation fixes the cusp \(\infty\), so it lies in \(H\).
This proves the second assertion from the first; the reverse implication
follows from \(\Gamma(N)\subseteq\Gamma_1(N)\).

The restriction to \(\Gamma_1(N)\) is the character induced by
\(\pi^*\iota(P_\lambda)\) on \(X_1(N)\). The construction by paths
above also gives a surjection from \(\Gamma_1(N)\) to the fundamental
group of its compact modular curve. Consequently the restricted character
is zero exactly when \(\pi^*\iota(P_\lambda)\) is trivial. By definition
this is the third condition. The zero character contributes the identity
in both groups.
\end{proof}

\begin{corollary}\label{cor:no2}
If \(E(\Q)[2]=0\), every nonzero mod-two elliptic period character
has noncongruence kernel.
\end{corollary}
\begin{proof}
The Shimura subgroup is the Cartier dual of a constant finite quotient
of the diamond group
\((\Z/N\Z)^\times/\{\pm1\}\) \cite{LO,Rib}. Its geometric two-torsion
consists of homomorphisms from that quotient to \(\mu_2\). All such
points are rational over \(\Q\). As \(\iota\) is an embedding over \(\Q\),
a nontrivial point of its intersection with \(\Sigma(N)[2]\) would give a
nontrivial point of \(E^\vee(\Q)[2]=E(\Q)[2]\), using the canonical
principal polarization of \(E\). Apply Theorem~\ref{th:criterion}.
\end{proof}

For an elliptic curve, the hypothesis in Corollary~\ref{cor:no2} is
equivalent to irreducibility of \(E[2]\) over \(\F_2\). Every stable line
in that two-dimensional representation has a single nonzero vector,
which the Galois group must fix. When rational two-torsion exists,
Theorem~\ref{th:criterion} still requires testing membership in the
Shimura subgroup.

\section{The primitive imaginary period}

Let \(f\) be the normalized weight-two newform of \(E\) and let
\(\omega_E\) be a N\'eron differential. Write
\(\varphi^*\omega_E=c_E\,2\pi i f(z)\,dz\), choosing its sign so that
\(c_E>0\). Let \(\Lambda_E\subset\C\) be the periods of \(\omega_E\)
on integral homology, and define \(\Omega_E^->0\) by
\(\Lambda_E\cap i\R=i\Omega_E^-\Z\). Put
\(c_\infty=\#\pi_0(E(\R))\). These conventions keep the primitive
anti-invariant period distinct from the generator of the projected
imaginary lattice; see also Wuthrich \cite{Wut}.

\begin{proposition}
The primitive imaginary period character on \(\Gamma_0(N)\) is
\[
 \chi^-(\gamma)=
 \frac{2c_E}{c_\infty\Omega_E^-}\,
 \operatorname{Im}\!\left(2\pi i\int_z^{\gamma z} f(u)\,du\right)
 \pmod 2 .
\]
It is nonzero.
\end{proposition}

\begin{proof}
In \(\C/\Lambda_E\), a class represented by \(v\) is real exactly when
\(2i\operatorname{Im}v\in\Lambda_E\).
Its possible imaginary coordinates therefore form
\((\Omega_E^-/2)\Z\). Modulo translations by \(\Lambda_E\), each such
horizontal line is one connected component (its image is a real circle).
It follows that
\[
 [(\Omega_E^-/2)\Z:\operatorname{Im}\Lambda_E]=c_\infty,
 \qquad
 \operatorname{Im}\Lambda_E=(c_\infty\Omega_E^-/2)\Z .
\]
Division by this last generator gives a surjective integral functional
on the elliptic period lattice. Pullback through \(\varphi_*\) and the
differential identity give the formula. Surjectivity gives its
nonzero reduction modulo two.
\end{proof}

For example, consider the isogeny class of
\(A:y^2=x^3-28x+1\). The cubic has no rational root: its only candidates
are \(1\) and \(-1\). Thus \(A[2]\) is irreducible. Take a
\(\Q\)-isogeny of least degree from \(A\) to the optimal elliptic curve in
its isogeny class. If its degree were even, its kernel would contain a
nonzero stable subgroup of \(A[2]\), hence all of \(A[2]\).
Factoring through multiplication by two would lower the degree. The
isogeny has odd degree and identifies the two-torsion representations.
Consequently each of the three nonzero period characters of this optimal
quotient has noncongruence kernel. This property concerns an unramified
modular cover. An application to the Birch and Swinnerton-Dyer rank
formula would still require a relation with the central order of the
complex \(L\)-function.

\begin{thebibliography}{9}
\bibitem{Woh} K. Wohlfahrt, \emph{An extension of F. Klein's level
concept}, Illinois J. Math. \textbf{8} (1964), 529--535.
\href{https://doi.org/10.1215/ijm/1256059574}{doi:10.1215/ijm/1256059574}.
\bibitem{LO} S. Ling and J. Oesterl\'e, \emph{The Shimura subgroup of
\(J_0(N)\)}, Ast\'erisque \textbf{196--197} (1991), 171--203.
\href{https://doi.org/10.24033/ast.75}{doi:10.24033/ast.75}.
\bibitem{Rib} K. A. Ribet, \emph{Congruence relations between modular forms},
\S4. \href{https://math.berkeley.edu/~ribet/Articles/icm.pdf}{Author's manuscript}.
\bibitem{Wut} C. Wuthrich, \emph{On the integrality of modular symbols and
Kato's Euler system for elliptic curves}, Doc. Math. \textbf{19} (2014).
\href{https://www.maths.nottingham.ac.uk/plp/pmzcw/download/int_kato.pdf}{Author's manuscript}.
\end{thebibliography}
\end{document}
